        本报告将证据分成四层：公司公告和产品页优先，其次是政策数据库和官方媒体，再次是券商公开PDF或转载，最后才是主题媒体梳理。所有关于WF6产能、收入占比、订单和扩产的判断都以公司公告或公司产品页为上限，市场传闻只能作为情绪指标。

        \begin{exhibitbox}[可入正文的高影响证据]
        \scriptsize
\resizebox{\exhibitwidth}{!}{
\begin{tabular}{llll}
\toprule
\textbf{结论} & \textbf{来源} & \textbf{置信度} & \textbf{处理} \\
\midrule
WF6用于CVD钨沉积和集成电路配线材料 & 昊华气体产品页 & 高 & 技术链条基础 \\
中船特气WF6现有产能2000吨/年、6N级 & 新华财经/证券时报 & 高 & 直接产能锚，订单仍需金额和平均售价 \\
中巨芯WF6产能600吨，暂无扩产计划 & 财联社/新浪财经 & 高 & 直接产能但有订单边界 \\
昊华科技2025年WF6收入占比0.13\% & 财联社/公告转述 & 高 & 避免过度放大单品 \\
华特气体无WF6合成产能，仅供应纯化WF6且占比较小 & 证券时报互动平台转述 & 高 & 下调WF6单品弹性 \\
WF6出口均价149.79美元/千克，同比+28.33\%、环比+203.83\% & 新华财经 & 中 & 只进情绪锚，不进合同平均售价 \\
钨价上涨来自供给、政策和战略需求共振 & 新华网/证券时报、新华社政策报道 & 中 & 资源端景气锚 \\
\bottomrule
\end{tabular}
}
\end{exhibitbox}

价格证据需要单独设门禁。钨价、WF6出口均价、海外报价通知和制冷剂配额都能解释为什么市场愿意给产业链更高关注度，但它们不等于公司实际成交价。尤其是WF6，公开出口均价只能说明行业供需紧张，不能证明中船特气、中巨芯-U或昊华科技拿到相同价格、相同毛利和相同客户结构。

\begin{exhibitbox}[公开价格/成本代理边界]
\footnotesize
\begin{tabularx}{\exhibitboxwidth}{L{1.0cm} L{1.7cm} L{1.8cm} X X X}
\toprule
\textbf{主题} & \textbf{代理项} & \textbf{数值} & \textbf{趋势} & \textbf{估值用途} & \textbf{仍缺字段} \\
\midrule
钨 & 65黑钨精矿 & 104.5万元/吨 & 较2026年初约+128\%，较2025年初约+633\% & 用于钨矿端周期景气、资源溢价和库存利润代理 & 公司矿端销售价格、库存成本、半导体高纯钨材客户订单仍需公告或年报验证 \\
钨 & APT & 150万元/吨 & 较2026年初超过+123\%，较2025年初约+615\% & 用于钨冶炼和粉末链价格传导代理 & 公司产品销售单价、加工费、存货计价和客户结构仍未完全披露 \\
钨 & 钨粉 & 2345元/千克 & 较2026年初超过+119\%，较2025年初约+644\% & 用于钨粉、硬质合金和WF6上游成本压力代理 & 无法替代高纯钨粉合同价、产品毛利或半导体客户订单 \\
钨 & 80钨铁 & 143万元/吨 & 较2026年初约+115\%，较2025年初约+562\% & 用于钢铁合金和高端制造需求侧价格验证 & 不直接对应A股公司高端钨材订单或电子材料价格 \\
电子特气 & 中国WF6出口均价 & 149.79美元/千克 & 同比+28.33\%，环比+203.83\% & 用于WF6行业供需紧张和出口价格情绪锚 & 公司合同平均售价、客户结构、销量、收入确认比例和产品毛利仍未披露 \\
电子特气 & 日韩WF6报价调整 & 报价上调70\%-90\% & 海外供给收缩叙事强化 & 只作为海外紧缺和国产替代情绪锚 & 不是国内上市公司合同平均售价，不能直接进入每股收益或目标价上修 \\
电子特气 & WF6钨粉成本占比 & 钨粉成本占60\%-70\% & 钨价上行会抬高WF6原料成本和议价压力 & 用于判断WF6涨价是否可能被成本吞噬 & 缺公司采购成本、价格传导条款和单品毛利率 \\
氟化工 & 2026三代制冷剂配额 & 三代制冷剂合计约79.78万吨，较2025年增加5963吨 & R32/R134a/R125配额和景气度仍是现金流主线 & 用于巨化、三美、永和等制冷剂现金流景气代理 & 公司单品价差、合同价格、库存和成本仍需季度财报验证 \\
\bottomrule
\end{tabularx}
\sourcenote{公开价格代理来自归档网页和AkShare宏观指数；行业价格不等于公司合同平均售价。}
\end{exhibitbox}

制冷剂配额也需要拆成两层看。行业级配额能解释板块景气，但不能告诉我们哪家公司有多少可验证的供给约束；生态环境部2026年度氢氟碳化物配额附件把这个缺口补到公司/实体层级，并能进一步算出全国份额。巨化体系全国生产配额份额约30.18\%，三美约9.64\%，永和约7.56\%；多氟多在该附件中没有检出同名配额行。因此，三美、巨化、永和的制冷剂基准现金流证据更扎实；多氟多仍应回到电子级HF、氟化盐和锂电材料本身，而不能套用制冷剂配额龙头逻辑。

\begin{exhibitbox}[公司级制冷剂配额证据]
\scriptsize
\begin{tabularx}{\exhibitboxwidth}{L{0.75cm} L{1.0cm} X L{1.55cm} L{1.75cm} X X}
\toprule
\textbf{代码} & \textbf{名称} & \textbf{映射实体} & \textbf{生产配额/全国份额} & \textbf{内用/进口配额} & \textbf{重点品种份额} & \textbf{估值用途/边界} \\
\midrule
600160 & 巨化股份 & 浙江衢化氟化学有限公司；浙江巨化股份有限公司电化厂；浙江兰溪巨化氟化学有限公司 & 240,826 / 30.18\% & 内用 120,716；进口 0 / 0.00\% & HFC-32 35.33\%；HFC-134a 28.66\%；HFC-125 29.86\%；HFC-143a 44.88\% & 补强制冷剂配额现金流、供给约束和基准业务估值；不直接上修电子材料第二曲线。；配额实体到上市公司合并口径仍需年报或公告核验；不得替代电子级HF、含氟电子材料客户、订单、单品价差、收入占比或毛利率。 \\
603379 & 三美股份 & 江苏三美化工有限公司；浙江三美化工股份有限公司 & 76,910 / 9.64\% & 内用 30,066；进口 8,451 / 0.14\% & HFC-134a 23.45\%；HFC-125 7.01\%；HFC-32 2.81\%；HFC-143a 15.90\% & 补强三美制冷剂主业现金流和配额景气的可验证性。；配额证明供给约束和主业现金流基础，不证明电子材料订单、合同价格、客户份额或单品毛利。 \\
605020 & 永和股份 & 内蒙古永和氟化工有限公司；金华永和氟化工有限公司；邵武永和金塘新材料有限公司；浙江永和制冷股份有限公司 & 60,281 / 7.56\% & 内用 25,524；进口 3,371 / 0.06\% & HFC-143a 28.57\%；HFC-134a 6.10\%；HFC-152a 31.59\%；HFC-125 5.37\% & 补强永和一体化氟化工和制冷剂配额主业基准。；制冷剂配额可进入基准现金流判断，但不能替代含氟电子材料客户、订单、价差、收入占比或毛利率。 \\
002407 & 多氟多 & 附件2未检出多氟多或同名HFC生产/进口配额实体 & 0 / 0.00\% & 内用 0；进口 0 / 0.00\% & 附件未检出 & 多氟多基准仍应看氟化盐、电子级HF、锂电材料和电子材料业务，不应套用制冷剂配额龙头逻辑。；未出现在本附件不等于公司没有其他氟化工业务；只说明本轮HFC制冷剂配额证据不能支持其制冷剂现金流估值。 \\
\bottomrule
\end{tabularx}
\sourcenote{生态环境部2026年度氢氟碳化物生产、进口配额核发表附件2；生产配额和内用生产配额单位为吨，进口配额单位为吨二氧化碳当量。全国份额以附件2附表1全部生产配额和附表2全部进口配额为分母聚合。实体到上市公司合并口径仍需年报或公告核验；该表只补强制冷剂基准现金流和相对地位，不替代含氟电子材料客户、订单、单品价差、收入占比或毛利率。}
\end{exhibitbox}

招股书证据进一步说明，价格缺口不是简单的公开渠道未检索到，而是发行人将六氟化钨和三氟化氮平均单价作为敏感商业信息处理。它能补强历史产销、价格降幅和毛利波动分析，但不能把公开出口均价倒填为中船特气当前合同平均售价；也不能把历史毛利率变动方向直接写成当前单品盈利。

\begin{exhibitbox}[中船特气招股书价格/毛利边界]
\scriptsize
\begin{tabularx}{\exhibitboxwidth}{L{1.5cm} X X X}
\toprule
\textbf{证据层} & \textbf{招股书披露} & \textbf{模型用途} & \textbf{估值边界} \\
\midrule
平均售价披露边界 & 招股说明书说明三氟化氮、六氟化钨平均单价按合计销售收入除以合计销售数量计算，并已申请豁免披露两种产品平均单价。 & 证明公司合同平均售价不是公开可得字段，公开出口均价、海外报价或行业报价不得替代公司销售均价。 & 只能作为P0硬缺口证明，不进入增长收入、增长每股收益或基本面目标价上修。 \\
历史产销可量化 & 招股说明书披露2019年、2020年、2021年、2022年上半年六氟化钨销量分别为269.19吨、533.75吨、838.27吨、550.15吨；产能利用率分别为33.62\%、75.83\%、58.44\%、56.80\%。 & 可支持中船特气六氟化钨直接性、历史产销兑现和容量监测。 & 销量和利用率只能约束产能兑现，不替代销售单价、收入确认比例或单品毛利。 \\
价格/成本/毛利敏感性 & 招股说明书披露六氟化钨平均单价2020年、2021年、2022年上半年分别下降8.27\%、5.76\%、5.16\%；2022年上半年毛利率下降7.50\%，原因是平均销售价格降低且单位成本上涨。 & 支持价格下行、成本上行和毛利率波动的风险折扣与监测阈值。 & 只能进入利润敏感性和风险折扣，不可倒推出当前合同平均售价或当前单品毛利。 \\
供需缺口历史基准 & 招股说明书按2022年至2025年国内六氟化钨需求1600吨、2400吨、2900吨、4500吨和国内产能2730吨、3530吨、3530吨、4030吨测算，2025年供需缺口为470吨。 & 可作为长期供需叙事和历史基准，说明六氟化钨需求增长和国产供给匹配问题。 & 2023年招股书预测必须用2026年公告、年报或客户订单刷新；不得替代当前订单、合同价格或收入确认。 \\
\bottomrule
\end{tabularx}
\sourcenote{中船特气招股说明书归档整理；该证据证明平均售价披露边界、历史产销和价格/毛利敏感性，不替代当前合同平均售价、收入确认比例或单品毛利。}
\end{exhibitbox}

因此，本报告对电子特气高成长业绩采取更硬的门禁：有历史销量和产能利用率，可以提高公司直接性和监测优先级；有历史平均价格降幅和毛利率变动，可以设置风险折扣；但只要当前合同平均售价、收入确认比例和单品毛利没有公告级披露，就不得把六氟化钨价格叙事转成增长每股收益。

\begin{exhibitbox}[公告/互动/研报补证结果]
\footnotesize
\begin{tabularx}{\exhibitboxwidth}{L{0.9cm} L{1.2cm} X X X}
\toprule
\textbf{代码} & \textbf{名称} & \textbf{已补证据} & \textbf{仍缺字段} & \textbf{估值门禁} \\
\midrule
688146 & 中船特气 & 2000吨/年、6N级；2025年六氟化钨采购合同约1.1904亿元；2024年三氟化氮销售合同约1.1508亿元；2025年HF/HCL/D2等11种电子气体产品采购合同约1.4822亿元；3383吨电子气体项目有公告/环评证据，但WF6单品新增拆分未由官方证实 & 合同平均售价、收入确认比例、客户份额、产品毛利、具名客户结构、2026年后持续订单和WF6单品新增产能拆分 & 历史主题订单金额样本可提高直接性和订单可见度，但不能把合同金额或项目总产能直接折成长期增长每股收益 \\
688549 & 中巨芯-U & 高纯WF6产能600吨；暂无扩产计划；无新增长期或大额实质性订单 & 具名客户、订单金额、价格、收入占比、产品毛利 & 直接产能存在，但估值只能保留市销率和观察信用，不给增长每股收益信用 \\
600378 & 昊华科技 & 2025年度WF6收入占公司营业收入0.13\%；官网产品页披露电子级WF6能力 & 订单金额、价格、客户结构、产品毛利；且单品占比低 & 按平台/分部估值观察，不得套用中船特气同等WF6弹性 \\
688268 & 华特气体 & 无WF6合成产能；仅按部分客户要求供应纯化WF6，业务占比较小 & 客户名称、订单金额、收入占比、产品毛利 & 电子特气平台价值保留，WF6单品只给小权重观察信用 \\
主题层 & WF6行业 & 公开报道显示出口均价和市场需求存在上行叙事，AI/HBM/3D NAND为需求解释 & 公司合同平均售价、季度销量、产品毛利和客户验证 & 行业价格只能解释主题热度，不能直接填入单公司收入桥 \\
\bottomrule
\end{tabularx}
\end{exhibitbox}
